    \documentclass[a4paper,12pt]{article}

\usepackage[utf8]{inputenc}
\usepackage[T1]{fontenc}
\usepackage[slovak]{babel}
\usepackage{graphicx}
\usepackage{amsmath}
\usepackage{hyperref}
\usepackage{geometry}
\usepackage{tikz}
\usetikzlibrary{arrows.meta, positioning}

\geometry{margin=2.5cm}

\title{Supervízne riadenie autonómnej skupiny dronov vo virtuálnej realite}
\author{Projektová dokumentácia}
\date{\today}

\begin{document}

%------------------------------------------------
% TITULNA STRANA
%------------------------------------------------
\begin{titlepage}
\centering

\vspace*{3cm}

{\LARGE \textbf{Supervízne riadenie autonómnych dronov pomocou virtuálnej reality}}

\vspace{1.5cm}

{\Large Projektová dokumentácia}

\vspace{2cm}

\textbf{Autori:} Bc. Marek Malina, Bc. Tadeáš Kapina, Bc. Šimon Maco, Bc. Miroslav Beňušík, Bc. Ľuboš Valkovič, Bc. Boris Šupák, Bc. Martin Brandobur, Bc. Ivan Vanca

\vfill

\today

\end{titlepage}

%------------------------------------------------
% OBSAH
%------------------------------------------------
\tableofcontents
\newpage

%------------------------------------------------
\section{Analýza zadania a požiadaviek}

Cieľom projektu je návrh a implementácia systému umožňujúceho supervízne riadenie skupiny autonómnych dronov pomocou virtuálnej reality (VR). 
Drony vykonávajú autonómne inšpekčné misie vo vnútornom priestore, pričom operátor má možnosť sledovať ich stav a v prípade potreby 
prevziať riadenie niektorého z dronov pomocou VR ovládačov.

Na rozdiel od jednoduchých systémov určených pre jeden dron musí navrhované riešenie podporovať viacero dronov súčasne. 
Každý dron vykonáva vlastnú autonómnu misiu, pričom operátor má možnosť sledovať ich činnosť a vybrať konkrétny dron, 
do ktorého riadenia môže zasiahnuť. Tento prístup umožňuje efektívne monitorovanie a koordináciu viacerých autonómnych 
robotických systémov v rovnakom prostredí.

Hlavnou požiadavkou systému je bezpečný prechod medzi autonómnym a manuálnym režimom riadenia. 
Operátor musí byť schopný zasiahnuť do riadenia bez destabilizácie letu a následne umožniť dronu 
vrátiť sa späť do autonómneho režimu.
\\

\textbf{Hlavné požiadavky systému:}

\begin{itemize}
    \item autonómne vykonávanie misií viacerými dronmi
    \item možnosť výberu konkrétneho drona zo zoznamu dostupných zariadení
    \item zobrazenie telemetrických dát a video streamu z jednotlivých dronov
    \item manuálne riadenie vybraného drona pomocou VR ovládačov
    \item bezpečný návrat do autonómneho režimu po manuálnom zásahu operátora
\end{itemize}

Systém musí okrem zabezpečenia požadovanej funkcionality spĺňať aj viaceré prevádzkové a technické požiadavky. 
Dôležitá je najmä spoľahlivá komunikácia medzi jednotlivými komponentmi systému, 
nízka latencia pri prenose riadiacich príkazov a implementácia bezpečnostných mechanizmov, 
ktoré minimalizujú riziko straty kontroly nad dronmi počas prevádzky.

%------------------------------------------------
\section{Návrh základnej architektúry systému}

Celý systém je navrhnutý ako distribuovaná architektúra pozostávajúca z troch hlavných častí:

\begin{itemize}
\item autonómny systém dronov
\item supervízna riadiaca vrstva
\item používateľské rozhranie vo virtuálnej realite
\end{itemize}

Autonómny systém každého drona zabezpečuje stabilizáciu letu, lokalizáciu a vykonávanie plánovanej misie. 
Supervízna vrstva koordinuje jednotlivé drony a umožňuje operátorovi sledovať ich stav. 
Používateľské rozhranie vo VR poskytuje operátorovi prehľad o všetkých aktívnych dronoch a 
umožňuje zasiahnuť do riadenia vybraného zariadenia.

Komunikácia medzi jednotlivými komponentmi systému je realizovaná pomocou middleware ROS~2, 
ktorý umožňuje modulárnu architektúru a jednoduchú integráciu jednotlivých komponentov.

\begin{figure}[h!]
\centering

\begin{tikzpicture}[
block/.style={draw, rectangle, rounded corners, align=center, minimum width=6cm, minimum height=1.4cm},
line/.style={-Latex, thick},
node distance=2cm
]

\node[block] (vr) {Používateľské rozhranie\\vo virtuálnej realite};

\node[block, below=of vr] (supervisor) {Supervízna riadiaca vrstva};

\node[block, below=of supervisor] (drone) {Autonómny systém dronov};

% komunikácia VR <-> supervisor
\draw[line] (vr) -- node[right]{riadiace príkazy} (supervisor);
\draw[line] (supervisor) -- node[left]{telemetria, stav systému} (vr);

% komunikácia supervisor <-> drony
\draw[line] (supervisor) -- node[right]{príkazy misie} (drone);
\draw[line] (drone) -- node[left]{telemetria, video stream} (supervisor);

\end{tikzpicture}

\caption{Bloková schéma architektúry navrhovaného systému}
\label{fig:system_architecture}

\end{figure}

Na obrázku \ref{fig:system_architecture} je znázornená základná architektúra navrhovaného systému. 
Používateľ komunikuje so systémom prostredníctvom rozhrania vo virtuálnej realite, ktoré umožňuje 
vizualizáciu telemetrických údajov a manuálne riadenie vybraného drona. Supervízna riadiaca vrstva 
zabezpečuje dohľad nad jednotlivými dronmi, výber aktívneho drona a prepínanie medzi autonómnym 
a manuálnym režimom riadenia. Autonómny systém dronov vykonáva samotnú misiu a zabezpečuje 
stabilizáciu letu, lokalizáciu a navigáciu jednotlivých dronov. Komunikácia medzi jednotlivými 
vrstvami systému zahŕňa prenos riadiacich príkazov smerom k dronom a prenos telemetrie a video 
streamu smerom k operátorovi.

%------------------------------------------------
\section{Návrh komunikačných rozhraní a ROS 2 uzlov}

Systém je implementovaný pomocou architektúry ROS~2 založenej na uzloch (nodes), ktoré komunikujú pomocou tém (topics), služieb (services) a akcií (actions).

Základné uzly systému sú:

\begin{itemize}
\item \textbf{Flight Interface Node} – komunikácia medzi ROS~2 a riadiacou jednotkou drona,
\item \textbf{State Estimation Node} – výpočet polohy a orientácie drona,
\item \textbf{Mission Executor Node} – vykonávanie plánovanej trajektórie,
\item \textbf{Supervisor Node} – riadenie režimov systému,
\item \textbf{VR Interface Node} – komunikácia medzi VR aplikáciou a ROS~2.
\end{itemize}

Medzi hlavné komunikačné témy patria napríklad:

\begin{itemize}
\item \texttt{/drone/state}
\item \texttt{/drone/pose}
\item \texttt{/drone/battery}
\item \texttt{/drone/cmd\_vel}
\item \texttt{/drone/mode}
\end{itemize}

Takáto architektúra umožňuje jednoduché rozšírenie systému o ďalšie drony alebo nové funkcionality.

%------------------------------------------------
\section{Návrh stavového automatu a handover mechanizmu}

Riadenie režimov systému je realizované pomocou stavového automatu (Finite State Machine – FSM). 
Tento automat zabezpečuje bezpečný prechod medzi jednotlivými režimami systému.

Základné stavy systému sú:

\begin{itemize}
\item \textbf{Idle} – systém čaká na spustenie misie,
\item \textbf{Autonomous Mission} – dron vykonáva autonómnu trajektóriu,
\item \textbf{Manual Control} – operátor riadi dron pomocou VR ovládačov,
\item \textbf{Hold} – dron udržiava stabilnú polohu,
\item \textbf{Landing} – bezpečné pristátie.
\end{itemize}

Pri prechode z autonómneho do manuálneho režimu systém najprv aktivuje stabilizačný režim 
(\textit{Hold}), čím sa zabezpečí plynulý prechod bez náhlych zmien trajektórie.

Podobne pri návrate z manuálneho riadenia do autonómneho režimu systém vypočíta najbližší bod 
trajektórie, od ktorého je možné bezpečne pokračovať v misii.

%------------------------------------------------
\section{Autonómna navigácia a vykonávanie misie}

Autonómna navigácia drona je založená na vykonávaní preddefinovanej trajektórie pozostávajúcej z 
niekoľkých waypointov. Každý waypoint predstavuje cieľovú polohu, ktorú má dron dosiahnuť.

Počas letu systém priebežne vyhodnocuje aktuálnu polohu drona a generuje riadiace príkazy, 
ktoré umožňujú sledovanie trajektórie.

Pre lokalizáciu vo vnútornom priestore je možné použiť rôzne metódy, napríklad:

\begin{itemize}
\item vizuálnu odometriu,
\item SLAM algoritmy,
\item externé lokalizačné systémy.
\end{itemize}

%------------------------------------------------
\section{VR používateľské rozhranie}

Používateľské rozhranie vo virtuálnej realite umožňuje operátorovi interagovať so systémom 
intuitívnym spôsobom. Operátor má k dispozícii vizuálne zobrazenie prostredia, telemetrické údaje 
a video stream z kamery drona.

Medzi hlavné funkcie VR rozhrania patria:

\begin{itemize}
\item výber dostupného drona zo zoznamu zariadení,
\item zobrazenie aktuálneho stavu drona,
\item vizualizácia video streamu,
\item manuálne riadenie pomocou VR ovládačov,
\item prepínanie medzi autonómnym a manuálnym režimom.
\end{itemize}

%------------------------------------------------
\section{Bezpečnostné mechanizmy a failsafe}

Bezpečnosť systému je kľúčovým aspektom návrhu. Systém musí byť schopný reagovať na rôzne 
poruchové stavy, ako napríklad strata komunikácie alebo neplatné riadiace príkazy.

Medzi hlavné bezpečnostné mechanizmy patria:

\begin{itemize}
\item watchdog kontrolujúci dostupnosť komunikácie,
\item automatický prechod do stabilizačného režimu pri strate spojenia,
\item obmedzenie maximálnej rýchlosti a zrýchlenia,
\item núdzové pristátie pri kritickom stave batérie.
\end{itemize}



%------------------------------------------------
\section{Implementácia na reálnom zariadení}

%------------------------------------------------
\section{Experimentálne overenie}

%------------------------------------------------
\section{Diskusia obmedzení a ďalší rozvoj}

\end{document}